% =============================================================
% Преамбула ENF для XeLaTeX.
%
% Цвета и макросы ролей подключаются из порождаемого файла
% enf-colors.tex — править его вручную нельзя, см. заголовок файла.
% =============================================================

% Цвета подключаются отдельным файлом через include-in-header в
% Build/defaults/pdf.yaml, а не командой \input: Pandoc копирует
% преамбулу во временный каталог, и относительный путь там не находится.

\usepackage{amsmath}
\usepackage{amssymb}
\usepackage{mdframed}
\usepackage{needspace}

% ---------- Коллауты ----------
% Lua-фильтр enf-callouts.lua превращает цитаты Obsidian в Div с классом
% callout-<тип>; Pandoc отображает такой Div окружением с тем же именем.
% Каждый тип получает не только цвет, но и словесный заголовок:
% ENF-COLOR-003 запрещает делать цвет единственным носителем смысла.

\newmdenv[
  linecolor=enfvariable, linewidth=2pt,
  topline=false, bottomline=false, rightline=false,
  backgroundcolor=white, skipabove=8pt, skipbelow=8pt,
  innerleftmargin=10pt, innertopmargin=6pt, innerbottommargin=6pt
]{enfcalloutdefinition}

\newmdenv[
  linecolor=enftarget, linewidth=2pt,
  topline=false, bottomline=false, rightline=false,
  skipabove=8pt, skipbelow=8pt,
  innerleftmargin=10pt, innertopmargin=6pt, innerbottommargin=6pt
]{enfcallouttheorem}

\newmdenv[
  linecolor=enfoperator, linewidth=2pt,
  topline=false, bottomline=false, rightline=false,
  skipabove=8pt, skipbelow=8pt,
  innerleftmargin=10pt, innertopmargin=6pt, innerbottommargin=6pt
]{enfcalloutproof}

\newmdenv[
  linecolor=enfparameter, linewidth=2pt,
  topline=false, bottomline=false, rightline=false,
  skipabove=8pt, skipbelow=8pt,
  innerleftmargin=10pt, innertopmargin=6pt, innerbottommargin=6pt
]{enfcalloutintuition}

\newmdenv[
  linecolor=enffunction, linewidth=2pt,
  topline=false, bottomline=false, rightline=false,
  skipabove=8pt, skipbelow=8pt,
  innerleftmargin=10pt, innertopmargin=6pt, innerbottommargin=6pt
]{enfcalloutexample}

\newmdenv[
  linecolor=enfneutral, linewidth=2pt,
  topline=false, bottomline=false, rightline=false,
  skipabove=8pt, skipbelow=8pt,
  innerleftmargin=10pt, innertopmargin=6pt, innerbottommargin=6pt
]{enfcalloutremark}

% Псевдонимы для типов, оформляемых одинаково с базовым
\newenvironment{enfcalloutlemma}{\begin{enfcallouttheorem}}{\end{enfcallouttheorem}}
\newenvironment{enfcalloutcorollary}{\begin{enfcallouttheorem}}{\end{enfcallouttheorem}}
\newenvironment{enfcalloutwarning}{\begin{enfcalloutremark}}{\end{enfcalloutremark}}
\newenvironment{enfcalloutnote}{\begin{enfcalloutremark}}{\end{enfcalloutremark}}

% ---------- Выключные формулы ----------
% Не разрывать формулу между страницами: разорванный вывод нечитаем.
\AtBeginDocument{\interdisplaylinepenalty=10000}

% ---------- Заголовок коллаута ----------
\newcommand{\enfcallouttitle}[1]{{\sffamily\bfseries #1}\par\vspace{2pt}}
